\section*{Chapter \ref{ch:in:crypto-native-firms} --- Crypto-Native Firms}

\begin{solution}{exo:in:crypto-native-firms:1}
Revenue: $-59.3\%$ (2022), $-2.7\%$ (2023), $+111.2\%$ (2024), $+9.4\%$ (2025). Headcount: $-24.3\%$ (2023), $+10.4\%$
(2024), $+31.3\%$ (2025); 2021's headcount is not in the filings read.
\end{solution}

\begin{solution}{exo:in:crypto-native-firms:2}
25\% (\$100\,000 at the grant price) at the cliff, then $1/48$ of the grant (\$8\,333) each month for 36 months.
\end{solution}

\begin{solution}{exo:in:crypto-native-firms:3}
Tax $0.4\times10\,000=\$4\,000$; proceeds \$5\,000; net \$1\,000.
\end{solution}

\begin{solution}{exo:in:crypto-native-firms:4}
With no drift the expected price stays constant, but the lognormal median falls as $e^{-\sigma^2t/2}$: a few paths rise a
lot and many fall. The mean is carried by the rare large outcomes; the typical outcome is below it.
\end{solution}

\begin{solution}{exo:in:crypto-native-firms:5}
\$17\,150; the remaining \$42\,850 ranks as a general unsecured claim.
\end{solution}

\begin{solution}{exo:in:crypto-native-firms:6}
2023 (headcount $-24\%$ as revenue fell only 3\%) and 2025 (headcount $+31\%$ as revenue rose 9\%): in both, headcount was
responding to the previous year's revenue, the collapse of 2022 and the doubling of 2024.
\end{solution}

\begin{solution}{exo:in:crypto-native-firms:7}
The median falls from \$106\,193 to \$76\,275, and the share of paths with a tranche whose tax exceeds its proceeds rises
from 76.2\% to 86.8\%: a longer lock-up leaves more time for the price to fall after tax is fixed.
\end{solution}

\begin{solution}{exo:in:crypto-native-firms:8}
Equal expectation is not equivalence: the grant's median is less than half the cash value, its outcomes range from
almost nothing to more than twice cash, and tax falls on value at vesting whatever happens later. A risk-averse employee
values the grant below its mean (chapter~13's certainty equivalent), and the firm's own survival is correlated with the
token.
\end{solution}

\begin{solution}{pb:in:crypto-native-firms:1}
\begin{enumerate}
\item Market makers and trading firms, funds, exchanges and brokers, protocol teams and foundations.
\item A non-profit that funds and coordinates an open-source protocol from a treasury partly in its token; one sells its
token periodically and builds fiat savings in bull markets to fund spending in bear markets.
\item Round-the-clock markets; self-custody; pay in tokens; unsettled regulatory status.
\item Revenue \$7.8, 3.2, 3.1, 6.6, 7.2 billion; employees (2022--2025) 4\,510, 3\,416, 3\,772, 4\,951.
\item June 2022, about 18\% (about 1\,100); January 2023, about 21\% of end-2022 headcount (about 950).
\item Pay in a crypto asset under a vesting schedule, often with a lock-up; a right attached to equity to receive future
tokens.
\item As wages (US) and as employment income subject to income tax and National Insurance (UK).
\item Usually at vesting, on the value then.
\item The employee has paid tax on value that has gone.
\item \$1\,000.
\item 25\% at twelve months, then $1/48$ monthly to month 48; each tranche locked six months.
\item Schedule; simulate prices; tax at vesting and proceeds after lock-up per tranche; distribution against cash.
\item \$106\,193; \$17\,467 and \$537\,810; \$235\,062.
\item \$240\,000 (\$400\,000 less 40\% tax); in 74.3\% of paths the grant is worth less.
\item 76.2\% at 80\% volatility; 2.4\% at 40\%.
\item Median \$106\,193, 10th--90th percentile \$17\,467 to \$537\,810 against \$240\,000 in cash; 76.2\% of paths with a
tranche taxed on more than it sells for.
\item Up to \$17\,150.
\item Joining at the top of a cycle puts her in the cohort hired on peak revenue and most exposed to the next cut.
\item Sell tranches as soon as allowed; hedge with futures or options on the token where they exist and her contract
allows; negotiate cash for part of the grant or tax at sale.
\item The cash offer is worth more to her unless she wants a lottery on the token and can bear losing most of it; if she
takes the grant, she should plan to sell on schedule and budget the tax in cash.
\end{enumerate}
\end{solution}

\begin{solution}{iq:in:crypto-native-firms:1}
Custody of client assets (key management, wallets), on-chain deposits and withdrawals with confirmations, round-the-clock
operation without a closing auction, and security against attackers who can steal irreversibly.
\iqlookfor{custody and irreversibility.}
\end{solution}

\begin{solution}{iq:in:crypto-native-firms:2}
The markets never close and move most when others sleep; a follow-the-sun desk hands the book between regions, with clear
handover of positions, limits and open incidents, and an on-call engineer per shift.
\iqlookfor{handover discipline.}
\end{solution}

\begin{solution}{iq:in:crypto-native-firms:3}
$e^{-\sigma^2/2}=e^{-0.32}=0.726$ of today's price.
\iqlookfor{the lognormal median.}
\end{solution}

\begin{solution}{iq:in:crypto-native-firms:4}
Short futures or perpetuals on the token, or buy puts, if they exist and are liquid; limits: the employment contract or
insider rules may forbid it, the market may be too thin, and funding costs may be high.
\iqlookfor{instruments, and the contractual and practical constraints.}
\end{solution}

\begin{solution}{iq:in:crypto-native-firms:5}
Expected withdrawal flow and its tail against the risk of holding keys online: enough in hot wallets to meet
withdrawals, the rest cold. Risks: theft from hot wallets, delays and operational errors in moving from cold storage.
\iqlookfor{liquidity against security.}
\end{solution}

\begin{solution}{iq:in:crypto-native-firms:6}
Correlate headcount changes with revenue changes at lags zero and one: here the lag-one relation is the stronger. With
four or five points, report it as a pattern, not a test, and look for other firms' data before generalising.
\iqlookfor{lags, and humility with five points.}
\end{solution}
